\documentclass[11pt,reqno]{amsart}
\usepackage[T1]{fontenc}
\usepackage{lmodern}
\usepackage[letterpaper,margin=1in]{geometry}
\usepackage{microtype}
\usepackage{amsmath,amssymb,mathtools,mathrsfs}
\usepackage[hidelinks,pdfusetitle]{hyperref}
\usepackage{enumitem}
\usepackage{needspace}
\numberwithin{equation}{section}
\newtheorem{theorem}{Theorem}[section]
\newtheorem{proposition}[theorem]{Proposition}
\newtheorem{lemma}[theorem]{Lemma}
\newtheorem{corollary}[theorem]{Corollary}
\theoremstyle{definition}
\newtheorem{definition}[theorem]{Definition}
\theoremstyle{remark}
\newtheorem{remark}[theorem]{Remark}
\newcommand{\R}{\mathbb R}
\newcommand{\Z}{\mathbb Z}
\newcommand{\N}{\mathbb N}
\newcommand{\E}{\mathbb E}
\newcommand{\1}{\mathbf 1}
\newcommand{\Conf}{\operatorname{Conf}}
\newcommand{\Var}{\operatorname{Var}}
\newcommand{\PV}{\operatorname{PV}}
\newcommand{\Ent}{\operatorname{Ent}}
\newcommand{\Sine}{\operatorname{Sine}}
\newcommand{\Id}{\mathrm{Id}}
\newcommand{\Leb}{\operatorname{Leb}}
\DeclareMathOperator{\Law}{Law}
\newcommand{\dd}{\,d}
\newcommand{\gap}{\mathfrak g}
\newcommand{\calB}{\mathcal B}
\newcommand{\calI}{\mathcal I}
\newcommand{\calE}{\mathcal E}
\newcommand{\calF}{\mathcal F}
\newcommand{\calG}{\mathcal G}
\newcommand{\DLR}{\mathcal{K}}
\newcommand{\fish}{\text{Fish}}
\newcommand{\norm}[1]{\left\lVert #1\right\rVert}
\newcommand{\abs}[1]{\left\lvert #1\right\rvert}
\DeclareMathOperator{\supp}{supp}
\DeclareMathOperator{\diver}{div}
\allowdisplaybreaks[2]
\setlength{\emergencystretch}{1.5em}
\setlist[enumerate]{leftmargin=*,itemsep=3pt,topsep=5pt}
\title[Uniqueness of the logarithmic DLR state]{Uniqueness of the stationary logarithmic DLR state}
\date{Research manuscript draft\\9 September 2026; revised and simplified 10 September 2026}
\subjclass[2020]{60G55, 60K35, 82B21, 60B20}
\keywords{Logarithmic gas, DLR equations, Sine beta process, hyperuniformity, Palm measure, entropy}

\begin{document}
\begin{proposition}[Entropy and Fisher-information bounds for gap averages]\label{prop:gap-entropy}
Fix $\beta>0$, an integer $r\ge1$, and $G\in C_c^2((0,\infty)^r)$.
There is a constant $C_G<\infty$ such that, for every integer
$K\ge2r+1$, every $T>0$, and every probability measure $\mu$ on
$\Omega_K$,
\begin{align}
|\mu(F_K)-\omega_{K,T}(F_K)|^2
&\le \frac{C_G}{K}S(\mu\mid\omega_{K,T}),\label{eq:entropy-gap}\\
|\mu(F_K)-\omega_{K,T}(F_K)|^2
&\le \frac{C_GT}{K}I(\mu\mid\omega_{K,T}).\label{eq:fisher-gap}
\end{align}
More precisely, choose $B<\infty$ so that
$\supp G\subset(0,B)^r$, and put
\[
M_0=\norm G_\infty,\qquad
M_1=\max_a\norm{\partial_aG}_\infty,\qquad
M_2=\max_{a,b}\norm{\partial_{ab}G}_\infty.
\]
The constant may be chosen to depend only on
$\beta,r,B,M_0,M_1,M_2$.
\end{proposition}


\begin{proof}
Write $\Phi=\Phi_{K,T}$ and $n=K-r\ge K/2$. For $v\in\R^K$,
extend $v$ by $v_0=v_{K+1}=0$ and put $d_i=v_{i+1}-v_i$.
Nearest-neighbor interactions and the two endpoint charges give
\begin{equation}\label{eq:hessian}
v^\top\Phi''(u)v
\ge\beta\sum_{i=0}^K\frac{d_i^2}{g_i(u)^2}
   +\frac1T|v|^2.
\end{equation}
Every remaining logarithmic interaction has nonnegative Hessian.
We use the stronger gap-dependent part of this lower bound to obtain
the factor $K^{-1}$; the uniform part alone would not give it.

\smallskip\noindent\textit{First gap derivatives.}
Regard $F_K$ as a function of the gaps and set
$q_j=\partial F_K/\partial g_j$. Each gap belongs to at most $r$
windows, so
\[
q_j=\frac1n\sum_{\substack{1\le i\le n\\i\le j\le i+r-1}}
  \partial_{j-i+1}G(g_i,\ldots,g_{i+r-1}),\qquad
|q_j|\le\frac{rM_1}{n}\le\frac{2rM_1}{K}.
\]
The endpoint coefficients $q_0,q_K$ vanish. Since the support of
each derivative is contained in $\supp G$, a nonzero $q_j$ requires
$g_j\le B$. Consequently
\[
\sum_{j=0}^K q_j^2g_j^2\le\frac{4r^2M_1^2B^2}{K}.
\]
Weighted Cauchy--Schwarz and \eqref{eq:hessian} now give
\[
|\nabla F_K\cdot v|^2
=\abs{\sum_jq_jd_j}^2
\le\left(\beta^{-1}\sum_jq_j^2g_j^2\right)
     \left(\beta\sum_j\frac{d_j^2}{g_j^2}\right)
\le\frac{A_1}{K}v^\top\Phi''v,
\qquad A_1=\frac{4r^2B^2M_1^2}{\beta}.
\]
Taking the supremum over $v$ in this quadratic-form duality yields
\begin{equation}\label{eq:gradient-metric}
\nabla F_K^\top(\Phi'')^{-1}\nabla F_K\le A_1/K.
\end{equation}

\smallskip\noindent\textit{Second gap derivatives and window overlap.}
Let $w_i=(g_i,\ldots,g_{i+r-1})$ and let $J_i$ be the indicator
that $w_i$ belongs to the support of $D^2G$. Then
\[
\begin{split}
|v^\top F_K''v|
&\le\frac{M_2}{n}\sum_{i=1}^n
  J_i\left(\sum_{a=0}^{r-1}|d_{i+a}|\right)^2\\
&\le\frac{rM_2}{n}\sum_{i=1}^nJ_i
  \sum_{a=0}^{r-1}d_{i+a}^2\\
&\le\frac{rM_2B^2}{n}\sum_{i=1}^n
  \sum_{a=0}^{r-1}\frac{d_{i+a}^2}{g_{i+a}^2}
\le\frac{r^2M_2B^2}{n}\sum_{j=0}^K\frac{d_j^2}{g_j^2}.
\end{split}
\]
In the third line, every gap in an active window is at most $B$;
in the last line, every gap is counted at most $r$ times. Thus
\begin{equation}\label{eq:tilt-hessian}
|v^\top F_K''v|
\le\frac{A_2}{K}
 v^\top(\beta H_K^{{\rm lat}\,\prime\prime})v,
\qquad A_2=\frac{2r^2B^2M_2}{\beta}.
\end{equation}

\smallskip\noindent\textit{Uniformity over the tilted laws.}
For $s\in\R$, define the exponentially tilted probability measure
\[
\omega_s(\dd u)=
\frac{e^{sF_K(u)}}{\omega_{K,T}(e^{sF_K})}\,\omega_{K,T}(\dd u),
\]
Set $c_G=[2(1+A_2)]^{-1}$. The same pointwise estimate
\eqref{eq:tilt-hessian}, with no dependence on $s,K,T$, gives
for every $|s|\le c_GK$
\[
(\Phi-sF_K)''
\ge\tfrac12\beta H_K^{{\rm lat}\,\prime\prime}+T^{-1}\Id_K
\ge\tfrac12\Phi''.
\]
The Brascamp--Lieb inequality in Lemma~\ref{lem:convex-approx},
followed by \eqref{eq:gradient-metric}, therefore gives
\[
\Var_{\omega_s}F_K
\le\E_{\omega_s}\bigl[
\nabla F_K^\top((\Phi-sF_K)'')^{-1}\nabla F_K\bigr]
\le\frac{2A_1}{K}
\qquad(|s|\le c_GK).
\]
This estimate first holds on each gap-truncated simplex and passes
to $\omega_s$ with the same constants, as proved in the appendix.
If $\psi(s)=\log\omega_{K,T}(e^{sF_K})$, then
$\psi''(s)=\Var_{\omega_s}F_K$. Integrating twice from zero yields
\begin{equation}\label{eq:mgf}
\log\E_{\omega_{K,T}}
\exp\{s(F_K-\omega_{K,T}(F_K))\}
\le A_1s^2/K,\qquad |s|\le c_GK.
\end{equation}

\smallskip\noindent\textit{Entropy and Fisher information.}
The entropy inequality, used with either sign, gives
\[
|\mu(F_K)-\omega_{K,T}(F_K)|
\le\frac{S(\mu\mid\omega_{K,T})}{s}+\frac{D s}{K},
\qquad 0<s\le c_GK,\qquad D=1+A_1.
\]
Abbreviate $S=S(\mu\mid\omega_{K,T})$ for the rest of this paragraph.
When $0<S\le Dc_G^2K$, choose $s=\sqrt{SK/D}$ to obtain the
square bound $4DS/K$; the case $S=0$ follows by $s\downarrow0$.
When $S>Dc_G^2K$, use
$|\mu(F_K)-\omega_{K,T}(F_K)|\le2M_0$ instead.
Thus \eqref{eq:entropy-gap} holds, for example, with any
\[
C_G\ge\max\{4D,\,4M_0^2/(Dc_G^2)\}.
\]
When $S=+\infty$, \eqref{eq:entropy-gap} is immediate.
Finally, the uniform convexity in \eqref{eq:hessian} and the
logarithmic Sobolev part of Lemma~\ref{lem:convex-approx} imply
\begin{equation}\label{eq:lsi}
S(\mu\mid\omega_{K,T})\le\frac T2 I(\mu\mid\omega_{K,T}),
\end{equation}
which proves \eqref{eq:fisher-gap} after enlarging the constant if
necessary. The static convexity inequalities are also used in
\cite[Sections~3 and 5]{BEY}. The appendix makes their use on the
singular simplex, including the weak Fisher-information domain,
explicit for every $\beta>0$.
\end{proof}
\section{New Pro}
\begin{proposition}[Entropy and Fisher-information bounds for gap averages]\label{prop:gap-entropy}
Fix $\beta>0$, an integer $r\ge1$, and $G\in C_c^2((0,\infty)^r)$.
There is a constant $C_G<\infty$ such that, for every integer
$K\ge2r+1$, every $T>0$, and every probability measure $\mu$ on
$\Omega_K$,
\begin{align}
|\mu(F_K)-\omega_{K,T}(F_K)|^2
&\le \frac{C_GT}{K}\fish(\mu\mid\omega_{K,T}).\label{eq:fisher-gap}
\end{align}
\end{proposition}
\begin{proof}
Applying Lemmas \ref{lem:entropicControl} and \ref{lem:varFormula} to $|\mu(F_K)-\omega_{K,T}(F_K)|$ gives
  $$
|\mu(F_K)-\omega_{K,T}(F_K)|
\le
\frac{1}{s}\left[
\operatorname{Ent}(\mu\mid\nu)
+\int_0^s \int_0^u \Var_{\omega_t}(f)\ dt\ du
\right]
$$
where $\omega_t$ is the Gibbs measure on $\Omega_K$ with with potential $\DLR_\Z-sF_k$. Now, Lemma \ref{lem:computatialEstimates} and (trouver ref) entails that Poincaré's inequality holds for $\omega_t$ leading to
$$
\Var_{\omega_t}(f)\leq \E_{\omega_t}\left[\|\nabla F_K \|_{((\DLR^{\tiny pen}_\Z-sF_k)'')^{-1}}\right]
$$


\end{proof}

\begin{lemma}
\label{lem:entropicControl}
    Let $\mu$ and $\nu$ be two probability measures. Then, for any measurable function $f\in L^1(\mu)\cap L^1(\nu)$ and any $s>0$, we have
    $$
|\mu(f)-\nu(f)|
\le
\frac{1}{s}\left[
\operatorname{Ent}(\mu\mid\nu)
+
\max\!\left\{
\log \nu\!\left(e^{s(f-\nu (f))}\right),
\log \nu\!\left(e^{-s(f-\nu (f))}\right)
\right\}
\right].
$$
\end{lemma}
\begin{proof}
    This is a direct consequence of Donsker-Varadhan inequality applied to $sf$ then to $-sf$.
\end{proof}
\begin{lemma}
    Let $\nu$ be a probability measures. Then, for any bounded measurable function $f$ and any $s>0$, we have
    $$
\log \nu\!\left(e^{s(f-\nu (f))}\right)=\int_0^s \int_0^u \Var_{\nu_t}(f)\ dt\ du
$$
where $\nu_s$ is the probability such that $\nu_s(dx) \propto e^{sf(x)}\ \nu(dx) $
\end{lemma}
\begin{proof}
It simply follows from differentiating two times w.r.t.\ $s$ and Lebesgue differentiation theorem.
\end{proof}
\begin{lemma}
    \label{lem:computatialEstimates}
    \begin{itemize}
        \item 
        $$\|\nabla F_K\|_{(\DLR^{\tiny pen}_\Z)''}\lesssim K^{-1}$$
        \item 
        $$F_K '' \lesssim \frac{\beta}{K} \DLR''_{\Z} $$
        \item 
        \[
(\DLR^{\tiny pen}_{\Z}-sF_K)''
\ge\tfrac12\beta H_K^{{\rm lat}\,\prime\prime}+T^{-1}\Id_K
\ge\tfrac12 \DLR^{\tiny pen}_{\Z}''.
\]
    \end{itemize}
\end{lemma}
\end{document}
